\documentclass[11pt]{article}
\usepackage{amsmath}
\usepackage{listings}
\usepackage{color}
\usepackage{graphicx}
\usepackage{geometry}
\geometry{margin=1in}

% R code style for listings
\definecolor{mygreen}{rgb}{0,0.6,0}
\definecolor{mygray}{rgb}{0.5,0.5,0.5}
\definecolor{mymauve}{rgb}{0.58,0,0.82}
\lstset{
	language=R,
	basicstyle=\footnotesize\ttfamily,
	numbers=left,
	numberstyle=\tiny\color{mygray},
	stepnumber=1,
	numbersep=5pt,
	backgroundcolor=\color{white},
	showspaces=false,
	showstringspaces=false,
	showtabs=false,
	frame=single,
	rulecolor=\color{black},
	tabsize=2,
	captionpos=b,
	breaklines=true,
	breakatwhitespace=false,
	title=\lstname,
	keywordstyle=\color{blue},
	commentstyle=\color{mygreen},
	stringstyle=\color{mymauve}
}

\title{Problem Set 1: Education and Political Economy}
\author{QiLiu 25340516}
\date{\today}

\begin{document}
	\maketitle
	
	%%%%%%%%%%%%%%%%%%%%%%%%%%%%%%%%%%%%%%%%%%%%%%%%%%%%%%%%%%%%%%%%
	\section*{Problem 1: Education}
	
	A school counselor was curious about the average IQ of students in her school and took a random sample of 25 students' IQ scores:
	
	\begin{lstlisting}[language=R,caption={R code for Problem 1}]
		# Load libraries (none needed here)
		rm(list=ls())
		detachAllPackages <- function() {
			basic.packages <- c(
			"package:stats", "package:graphics", "package:grDevices", 
			"package:utils", "package:datasets", "package:methods", "package:base"
			)
			package.list <- search()[ifelse(unlist(gregexpr("package:", search()))==1, TRUE, FALSE)]
			package.list <- setdiff(package.list, basic.packages)
			if (length(package.list) > 0) {
				for (package in package.list) detach(package, character.only=TRUE)
			}
		}
		detachAllPackages()
		y <- c(
		105, 69, 86, 100, 82, 111, 104, 110, 87, 108,
		87, 90, 94, 113, 112, 98, 80, 97, 95, 111,
		114, 89, 95, 126, 98
		)
		# 1. 90% Confidence Interval
		n <- length(y)
		sample_mean <- mean(y)
		sample_sd <- sd(y)
		alpha <- 0.10    # 90% CI
		t_crit <- qt(1 - alpha/2, df = n - 1)
		se <- sample_sd / sqrt(n)
		ci_lower <- sample_mean - t_crit * se
		ci_upper <- sample_mean + t_crit * se
		cat("Sample mean:", sample_mean, "\n")
		cat("Sample SD:", sample_sd, "\n")
		cat("90% CI for population mean: [", ci_lower, ",", ci_upper, "]\n")
		# 2. One-sample t-test (test if mean > 100)
		t_test_result <- t.test(y, mu = 100, alternative = "greater", conf.level = 0.95)
		print(t_test_result)
	\end{lstlisting}
	
	\subsection*{Part 1: 90\% Confidence Interval for the Mean}
	
	The sample mean and standard deviation, calculated from the data, are:
	\begin{itemize}
		\item $\bar{y} = 98.44$
		\item $s = 13.09$
	\end{itemize}
	
	The 90\% confidence interval is given by:
	\[
	\left[ \bar{y} - t_{\alpha/2, n-1} \frac{s}{\sqrt{n}},\, \bar{y} + t_{\alpha/2, n-1} \frac{s}{\sqrt{n}} \right] =
	[93.96,\ 102.92]
	\]
	Therefore, we are 90\% confident that the true mean IQ score of students in this school is between 93.96 and 102.92.
	
	\subsection*{Part 2: Hypothesis Test}
	
	We test:
	\[
	\begin{aligned}
		H_0 &:~ \mu = 100, \\
		H_A &:~ \mu > 100
	\end{aligned}
	\]
	R output:
	\begin{verbatim}
		One Sample t-test
		data:  y
		t = -0.59574, df = 24, p-value = 0.7215
		alternative hypothesis: true mean is greater than 100
		95 percent confidence interval:
		93.95993      Inf
		sample estimates:
		mean of x 
		98.44 
	\end{verbatim}
	
	Since $p$-value $= 0.7215 \gg 0.05$, we \textbf{fail to reject} the null hypothesis. There is insufficient evidence to conclude that the average IQ at this school exceeds the national average of 100.
	
	\section*{Summary (Q1)}
	\begin{itemize}
		\item 90\% confidence interval: $\left[93.96,\ 102.92\right]$
		\item $p$-value = 0.72 ($t = -0.60$). The school's mean IQ is \textbf{not significantly higher} than the national average 100.
	\end{itemize}
	
	%%%%%%%%%%%%%%%%%%%%%%%%%%%%%%%%%%%%%%%%%%%%%%%%%%%%%%%%%%%%%%%
	\section*{Problem 2: Political Economy}
	
	Researchers are curious about what affects the amount communities spend on homelessness. The dataset has 50 US states (see code for variable descriptions).
	
	\subsection*{1. Preparation and Data Import}
	
	\begin{lstlisting}[language=R,caption={R code for data import}]
		rm(list=ls())
		detachAllPackages <- function() {
			basic.packages <- c("package:stats", "package:graphics", "package:grDevices", "package:utils",
			"package:datasets", "package:methods", "package:base")
			package.list <- search()[ifelse(unlist(gregexpr("package:", search()))==1, TRUE, FALSE)]
			package.list <- setdiff(package.list, basic.packages)
			if (length(package.list)>0) for (package in package.list) detach(package, character.only=TRUE)
		}
		detachAllPackages()
		# Install/load ggplot2
		pkgTest <- function(pkg){
			new.pkg <- pkg[!(pkg %in% installed.packages()[,"Package"])]
			if (length(new.pkg)){ install.packages(new.pkg, dependencies = TRUE)}
			sapply(pkg, require, character.only = TRUE)
		}
		pkgTest(c("ggplot2"))
		# Data import (demo - fill in with actual data as needed)
	    expenditure <- read.table("https://raw.githubusercontent.com/ASDS-TCD/StatsI_2025/main/datasets/expenditure.txt", header=T)
		str(expenditure); summary(expenditure); head(expenditure)
	\end{lstlisting}
	
	\textbf{Variables:}
	\begin{itemize}
		\item $Y$: Per capita expenditure on shelters/housing assistance
		\item $X_1$: Per capita income
		\item $X_2$: Residents per 100,000 "financially insecure"
		\item $X_3$: Per 1,000 living in urban areas
		\item Region: 1=Northeast, 2=North Central, 3=South, 4=West
	\end{itemize}
	
	\subsection*{2. Exploratory Data Analysis}
	
	\begin{lstlisting}[language=R,caption={Scatterplot matrix and correlation}]
		pairs(expenditure[,c("Y","X1","X2","X3")],
		main="Scatterplot Matrix: Y, X1, X2, X3")
		cor_mat <- cor(expenditure[,c("Y","X1","X2","X3")])
		print(round(cor_mat, 2))
	\end{lstlisting}
	
	\textbf{Figure 1: Scatterplot Matrix}\par
	\includegraphics[width=0.7\textwidth]{plot1.png}
	
	\textbf{Correlation output:}
	\begin{verbatim}
		Y   X1   X2   X3
		Y  1.00 0.53 0.45 0.46
		X1 0.53 1.00 0.21 0.60
		X2 0.45 0.21 1.00 0.22
		X3 0.46 0.60 0.22 1.00
	\end{verbatim}
	
	\subsection*{3. Housing Expenditure by Region}
	
	\begin{lstlisting}[language=R,caption={Boxplot for Y by Region}]
		region.lab <- c("Northeast","North Central","South","West")
		boxplot(Y ~ Region, data = expenditure, names = region.lab,
		ylab = "Per Capita Housing Assistance Expenditure",
		xlab = "Region",
		main = "Per Capita Expenditure by US Region",
		col = c("skyblue","lightgreen","orange","red"))
		region.means <- tapply(expenditure$Y, expenditure$Region, mean)
		print(region.means)
	\end{lstlisting}
	
	\textbf{Figure 2: Boxplot of Y by Region}\par
	\includegraphics[width=0.7\textwidth]{plot2.png}
	
\textbf{Region means:}

\begin{center}
	\begin{tabular}{cccc}
		1 & 2 & 3 & 4 \\
		79.44 & 83.92 & 69.19 & 88.31 \\
	\end{tabular}
\end{center}
	\textit{Interpretation:} West region (Region 4) has the highest mean per capita expenditure.
	
	\subsection*{4. Expenditure vs Personal Income}
	
	\begin{lstlisting}[language=R,caption={Scatterplot and regression line}]
		plot(expenditure$X1, expenditure$Y,
		xlab = "Per Capita Personal Income",
		ylab = "Per Capita Housing Assistance Expenditure",
		main = "Housing Assistance Expenditure vs Personal Income")
		abline(lm(Y ~ X1, data = expenditure), col="red", lwd=2)
	\end{lstlisting}
	
	\textbf{Figure 3: Expenditure vs Income}\par
	\includegraphics[width=0.7\textwidth]{plot3.png}
	
	\textit{Interpretation:} Higher per capita income tends to associate with higher housing assistance expenditure.
	
	\subsection*{5. Expenditure vs Income by Region}
	
	\begin{lstlisting}[language=R,caption={Scatterplot by region with ggplot2}]
		library(ggplot2)
		ggplot(expenditure, aes(x = X1, y = Y, color = factor(Region), shape = factor(Region))) +
		geom_point(size = 3) +
		labs(
		title = "Housing Assistance Expenditure vs Personal Income by Region",
		x = "Per Capita Personal Income",
		y = "Per Capita Housing Assistance Expenditure",
		color = "Region", shape = "Region"
		) +
		scale_color_manual(values = c("steelblue","forestgreen","orange","red"),
		labels = region.lab) +
		scale_shape_manual(values = 1:4, labels = region.lab) +
		theme_minimal()
	\end{lstlisting}
	
	\textbf{Figure 4: Expenditure vs Income (by Region)}\par
	\includegraphics[width=0.7\textwidth]{plot4.png}
	
	\textit{Interpretation:} The West (red) tends to spend more, even at similar income levels, indicating regional effects.
	
	\section*{Summary (Q2)}
	\begin{itemize}
		\item West region has the highest average per capita expenditure.
		\item Expenditure rises with income, but region is an independent factor.
		\item For similar incomes, Western states still spend more than others.
	\end{itemize}
	
\end{document}
